\section{A certified isolated boundary mode}
\label{boundary:sec:isolated-mode:rr}
\label{boundary:sec:certified-residue:rr}

A rational trial function, a full-operator negative-window exclusion,
and a bounded residue functional give the following joint certificate.

\begin{theorem}[A unique boundary mode with nonzero tail]
\label{v12:prop:certified-boundary-eigenvalue}
\label{boundary:thm:negative-exclusion:rr}
\label{boundary:thm:certified-residue:rr}
The phase-zero physical operator $B_0$ has no spectrum in
$[-953/2000,-713/2000]=[-0.4765,-0.3565]$ and exactly one eigenvalue
$\alpha_0$ in $[713/2000,953/2000]$. It is simple and lies in
\[
 (81/200,107/250)=(0.405,0.428).
\]
For every normalized physical eigenfunction at $\alpha_0$, the amplitude
in \eqref{v13:eq:physical-amplitude} satisfies $|\eta|>1/100$.
Consequently, along $\mathcal K_n$ with $n\in\mathbb N$ and $L=2n\to\infty$,
the even eigenvalue is eventually smaller than the odd one, with leading difference
$2|c_{\alpha_0}||\eta|^2L^{-\gamma(\alpha_0)}+o(L^{-\gamma(\alpha_0)})$.
\end{theorem}

Sections~\ref{boundary:app:step-parametrix:rr}
and~\ref{boundary:app:residue:rr} prove the three assertions in order:
a full-space residual gives existence, a global parametrix bound and
signed alternation give uniqueness, and physical spectral projection
transfers a certified residue value to the eigenfunction. All rounding
errors, omitted cells and spectral complements are included. The
nonvanishing statement concerns this mode; it imposes no condition on
other phases or boundary levels.
